\documentclass[a4paper,12pt]{article}
\usepackage[utf8]{inputenc}
\usepackage[T1]{fontenc}
\usepackage[ngerman]{babel}
\usepackage{amsmath}
\usepackage{amssymb}
\usepackage{enumitem}
\usepackage{geometry}
\geometry{a4paper, margin=2.5cm}

\title{Einsendeaufgaben -- Lektion 4\\[0.5em]
\large Modul 61111: Mathematische Grundlagen}
\author{}
\date{}

\begin{document}
\maketitle

\section*{Aufgabe 4.6}

Sei $a \in A_1 \cup A_2$ und $b \in B_1 \cap B_2$ beliebig.

\begin{enumerate}[label=\arabic*., start=2]

    \item zu zeigen: $(A_1 \cup A_2) \cup (B_1 \cap B_2) = \mathbb{R}$
    \begin{align*}
        (A_1 \cup A_2) \cup (B_1 \cap B_2)
        &= (A_1 \cup A_2 \cup B_1) \cap (A_1 \cup A_2 \cup B_2) \\
        &= (A_2 \cup \mathbb{R}) \cap (A_1 \cup \mathbb{R}) = \mathbb{R}
    \end{align*}

    \item $\forall a \in A_1 \cup A_2\; \forall b \in B_1 \cap B_2 : a \leq b$

    Sei $a \in A_1 \cup A_2$ und $b \in B_1 \cap B_2$ beliebig.

    \underline{Fall 1: $a \in A_1$}

    Da $b \in B_1$ folgt $a \leq b$.

    \underline{Fall 2: $a \in A_2$}

    Da $b \in B_2$ folgt $a \leq b$.

\end{enumerate}

Trennungszahl ist gleich $(A_1 | B_1)$ wenn $A_1 \supseteq A_2$, ansonsten andersrum.

\end{document}
